\documentclass[10pt,a4paper]{amsart}
\usepackage[margin=19mm]{geometry}
\usepackage{amsmath,amssymb,mathtools}
\usepackage[scr=rsfs,cal=euler]{mathalfa}
\usepackage{xfrac}
\usepackage[unicode,hidelinks,bookmarks=false]{hyperref}
\newcommand{\ie}{i.e.\ }
\newcommand{\cf}{cf.\ }
\newcommand{\resp}{respectively\ }
\newcommand{\Subsection}{\subsection}
\providecommand{\nobreakdash}{\nobreak\hbox{-}}
\setlength{\emergencystretch}{2em}
\multlinegap0pt
\usepackage{dsfont}
\usepackage{amssymb}
\usepackage{tikz}
\usepackage{float}
\usepackage[shortlabels]{enumitem}
\usepackage{algorithm}\usepackage{algpseudocode}
\usepackage{booktabs}
\usepackage{multirow}
\usepackage{bbm}
\usepackage{xcolor}
\newtheorem{problem}{Problem}
\newtheorem{theorem}{Theorem}
\newcommand{\calT}{{\mathcal T}}
\title{\Large \bfseries Temporally Flexible Transport Scheduling on Networks with Departure-Arrival Constriction and Nodal Capacity Limits}
\author{Anqi Dong, Karl H. Johansson, Johan Karlsson}
\date{Source: arXiv:2602.14652v1 (CC BY 4.0)}
\begin{document}
\maketitle

\noindent\emph{Source and licence.} \Large \bfseries Temporally Flexible Transport Scheduling on Networks with Departure-Arrival Constriction and Nodal Capacity Limits. Version arXiv:2602.14652v1 (CC BY 4.0), distributed by arXiv under the Creative Commons Attribution 4.0 International licence, http://creativecommons.org/licenses/by/4.0/, as recorded in the arXiv licence metadata for this exact version and shown on the abstract page https://arxiv.org/abs/2602.14652v1. Full licence notice: \texttt{licenses/arxiv-2602.14652-CC-BY-4.0.txt}.\par\smallskip

\noindent\textit{Editorial scope.} Counted unit: the article's theorem theorem (source0.tex lines 754--756). The article's complete original proof of it is delivered with it (source0.tex lines 759--765, one \texttt{proof} environment of 7 lines). No mathematics is written, repaired or re-derived: the statement and the proof are the article's own lines, in the article's own notation, and no retained line is substituted or re-flowed.  Retained uncounted as the source-local context the proof needs: lines 727--742.  Nothing was omitted from any retained span, so the package carries no omission record.  No retained line contains a citation, so the wrapper carries no bibliography and no bibliography entry.  No retained line contains a figure environment, an image-inclusion command or drawing code, so no image is delivered and the package declares no asset dependency.  Editorial formatting only, in addition: an AMS \texttt{amsart} preamble that restates the article's own declarations and macros used by the retained lines, namely the environments \texttt{problem}, \texttt{theorem} on separate counters, together with the standard packages those lines need.  The retained environments print in this document as Problem 1, Theorem 1 in order of appearance, whereas the article prints the same items with the numbering of its own full document.  The complete native source of the article as distributed in the arXiv e-print for this exact version is bundled unchanged as \texttt{sources/194733/source0.tex}.
\begin{problem}[\bf \textit{Coupled DA at single node}]\label{prob:prob2}
Given a coupled DA constraint represented by the probability measure $\mu^{0,2}$ and a flow-rate bound $r > 0$, determine a probability measure $\pi(t_0,t_1,t_2)$ on $\mathcal M_+([0, t_f]^3)$ that minimizes
\begin{subequations}
\begin{align}\label{eq:cost-function-3}
\int_{t_0,t_1,t_2} \Bigg( \frac{w(v_0,v_1)}{t_1-t_0} + \frac{w(v_1,v_2)}{t_2-t_1} \Bigg) \pi(dt_0,dt_1,dt_2)
\end{align}
subject to the coupled DA distribution
\begin{align}
\int_{t_1} \pi(dt_0,dt_1,dt_2) = \mu^{0,2}(dt_0,dt_2),
\end{align}
and flow-rate constraint $\int_{t_0,t_\calT} \pi(dt_0,dt_1,dt_2) \leq r dt_1$.
% \begin{align}
% \iint_{t_0,t_\calT} \pi(dt_0,dt_1,dt_\calT) \leq r dt_1.
% \end{align}
\end{subequations}
\end{problem}

\begin{theorem}[\bf \textit{Existence}]
For sufficiently large flow-rate bound $r$, Problem~\ref{prob:prob2} admits a minimizing solution $\pi$.
\end{theorem}

\begin{proof}
The space of admissible measures $\pi$ in Problem  \ref{prob:prob2} reads
\begin{align*}
\Pi = \Big\{&\pi \in \mathcal{M}_+\big([0, t_f]^3\big) \Big | \int_{t_1} \pi = \mu^{0,2}(dt_0,dt_2),\ \iint_{t_0, t_2} \pi:= \sigma^1(t_1) \leq r , dt_1 \Big\}
\end{align*}
is non-empty. It is compact in the weak topology as it is tight and closed for the narrow convergence. As its cost function is lower semi-continuous, we have the existence of a minimizer. 
\end{proof}

\end{document}
